\chapter{अवकल समीकरण}\label{ch:g12:diffeq}

\omterm{def:g12:diffeq:ode}{अवकल समीकरण} किसी \omterm{def:g11:func:function}{फलन} को उसके \omterm{def:g12:deriv:derivative}{अवकलजों} से जोड़ता है। भौतिकी, रसायन, जीवविज्ञान
और अर्थशास्त्र अपने नियम इसी रूप में व्यक्त करते हैं: किसी वस्तु के ठंडा होने
की दर, किसी रेडियोधर्मी नाभिक का क्षय, किसी समष्टि की वृद्धि --- ये सब $y'$ के
बारे में कथन हैं। यह अध्याय अचर गुणांकों वाले प्रथम कोटि के रैखिक \omterm{def:g10:algebra:equation}{समीकरण} पूरी
तरह हल कर देता है।

\section{समीकरण \texorpdfstring{$y' = ay$}{y' = ay}}

\begin{definition}[अवकल समीकरण]\label{def:g12:diffeq:ode}
\emph{अवकल समीकरण}\index{अवकल समीकरण} वह \omterm{def:g10:algebra:equation}{समीकरण} है जिसका अज्ञात कोई \omterm{def:g11:func:function}{फलन} $y$
हो और जिसमें $y$ तथा उसके \omterm{def:g12:deriv:derivative}{अवकलज} आते हों। किसी \omterm{def:g10:numbers:interval}{अंतराल} $I$ पर \emph{हल} वह
\omterm{def:g12:deriv:derivative}{अवकलनीय} \omterm{def:g11:func:function}{फलन} है जो $I$ के हर बिंदु पर \omterm{def:g10:algebra:equation}{समीकरण} को संतुष्ट करे।
\end{definition}

\begin{theorem}[$y' = ay$ के हल]\label{thm:g12:diffeq:homogeneous}
मान लीजिए $a \in \R$ है। \omterm{def:g10:algebra:equation}{समीकरण} $y' = ay$ के $\R$ पर हल ठीक ये \omterm{def:g11:func:function}{फलन} हैं
\[
y(x) = C\,\eu^{ax}, \qquad C \in \R .
\]
हर युग्म $(x_0, y_0)$ के लिए $y(x_0) = y_0$ वाला अद्वितीय हल होता है।
\end{theorem}

\begin{proof}
हर $y = C\eu^{ax}$ $y' = Ca\,\eu^{ax} = ay$ को संतुष्ट करता है। इसके उलट, मान लीजिए $y$ कोई हल है और $z(x) = y(x)\,\eu^{-ax}$
रखिए। तब
\[
z'(x) = y'(x)\,\eu^{-ax} - a\,y(x)\,\eu^{-ax}
= \bigl(y'(x) - a y(x)\bigr)\eu^{-ax} = 0,
\]
इसलिए $z$ अचर है, मान लीजिए $C$, और $y(x) = C\eu^{ax}$। आरंभिक शर्त $y(x_0) = y_0$ $C = y_0 \eu^{-a x_0}$ को
अद्वितीय रूप से तय कर देती है।
\end{proof}

\begin{remark}
इस तरह \omterm{thm:g12:exp:existence}{चरघातांकी फलन} का \emph{अभिलक्षण} सभी \omterm{def:g12:diffeq:ode}{अवकल समीकरणों} में सबसे सरल
\omterm{def:g10:algebra:equation}{समीकरण} से मिलता है: आकार के समानुपाती वृद्धि। इसीलिए वह प्रकृति में हर जगह
प्रकट होता है।
\end{remark}

\begin{omfigure}
\begin{tikzpicture}[scale=1.05]
  \foreach \x in {-2.8,-2.4,...,2.9}
    \foreach \y in {-1.8,-1.4,...,1.9}
      \draw[gray!60, thin] (\x,\y) ++({atan(\y)}:0.11)
        -- ++({atan(\y)+180}:0.22);
  \draw[-Stealth] (-3.2,0) -- (3.4,0) node[below, font=\small] {$x$};
  \draw[-Stealth] (0,-2.2) -- (0,2.4) node[left, font=\small] {$y$};
  \draw[omDef, thick] plot[domain=-3:0.69, samples=60] (\x, {exp(\x)});
  \draw[omThm, thick] plot[domain=-3:1.89, samples=60] (\x, {0.3*exp(\x)});
  \draw[omProp, thick] plot[domain=-3:1.38, samples=60] (\x, {-0.5*exp(\x)});
\end{tikzpicture}

{\small \omterm{def:g10:algebra:equation}{समीकरण} $y' = y$ समतल के हर बिंदु पर एक \omterm{def:g10:reffunc:affine}{प्रवणता} निर्धारित कर देता है
(धूसर रेखाखंड)। हल $C\eu^{x}$ --- यहाँ $C = 1$ (नीला), $C = 0.3$ (लाल), $C = -0.5$ (नारंगी) ---
ठीक वही वक्र हैं जो प्रवणताओं के इस क्षेत्र का अनुसरण करते हैं।}
\end{omfigure}

\begin{example}[रेडियोधर्मी क्षय]\label{ex:g12:diffeq:decay}
कोई रेडियोधर्मी राशि $\lambda > 0$ के साथ $N' = -\lambda N$ को संतुष्ट करती है, इसलिए $N(t) = N_0\,\eu^{-\lambda t}$;
अर्धायु के लिए \cref{ex:g12:exp:halflife} देखिए।
\end{example}

\section{समीकरण \texorpdfstring{$y' = ay + b$}{y' = ay + b}}

\begin{theorem}[$y' = ay + b$ के हल]\label{thm:g12:diffeq:affine}
मान लीजिए $a \neq 0$ और $b \in \R$ हैं। $y' = ay + b$ के $\R$ पर हल ठीक ये \omterm{def:g11:func:function}{फलन} हैं
\[
y(x) = C\,\eu^{ax} - \frac{b}{a}, \qquad C \in \R .
\]
अचर \omterm{def:g11:func:function}{फलन} $y_p = -\frac ba$ \emph{साम्य हल} है। हर $(x_0, y_0)$ के लिए $y(x_0) = y_0$ वाला अद्वितीय हल है;
और यदि $a < 0$ हो, तो $x \to +\infty$ पर हर हल साम्य $-\frac ba$ की ओर जाता है।
\end{theorem}

\begin{proof}
अचर $y_p = -\frac{b}{a}$ $y_p' = 0 = a y_p + b$ को संतुष्ट करता है। अब $y$ हल है यदि और केवल यदि
\[
(y - y_p)' = y' = ay + b = a(y - y_p) + \underbrace{a y_p + b}_{=\,0}
= a (y - y_p),
\]
अर्थात् यदि और केवल यदि $z = y - y_p$ $z' = az$ को हल करे। \cref{thm:g12:diffeq:homogeneous} से $z = C\eu^{ax}$ मिलता है, जिससे
सूत्र, अस्तित्व और अद्वितीयता निकल आते हैं। यदि $a < 0$ हो, तो $x \to +\infty$ पर $\eu^{ax} \to 0$,
इसलिए $y(x) \to -\frac ba$।
\end{proof}

\begin{method}[आरंभिक शर्त के साथ $y' = ay + b$ हल करना]\label{met:g12:diffeq:solve}
\begin{enumerate}
  \item साम्य हल $y_p = -\frac{b}{a}$ ज्ञात कीजिए ($y' = 0$ हल कीजिए)।
  \item व्यापक हल $y = C\eu^{ax} + y_p$ लिखिए।
  \item आरंभिक शर्त से $C$ ज्ञात कीजिए।
  \item दीर्घ काल के व्यवहार की जाँच भौतिक अंतर्ज्ञान के सामने कीजिए (क्या
        हल साम्य पर \omterm{def:g12:seq:limit}{अभिसरित} होता है?)।
\end{enumerate}
यही रणनीति --- \emph{विशेष हल + समघात \omterm{def:g10:algebra:equation}{समीकरण} का व्यापक हल} --- $y' = ay + f(x)$ तक भी
बढ़ती है: \cref{prop:g12:diffeq:general} देखिए।
\end{method}

\begin{example}[न्यूटन का शीतलन नियम]\label{ex:g12:diffeq:cooling}
$80\,^\circ$C की कॉफ़ी का प्याला $20\,^\circ$C के कमरे में रखा है। न्यूटन का नियम कहता है कि
ताप $T$ किसी $k > 0$ के लिए $T' = -k(T - 20)$ को संतुष्ट करता है, \emph{अर्थात्} $T' = -kT + 20k$। साम्य
$20$ है, और $T(t) = 20 + 60\,\eu^{-kt}$: कॉफ़ी चरघातांकी तेज़ी से कमरे के ताप तक ठंडी हो जाती है।
\end{example}

\begin{omfigure}
\begin{tikzpicture}
\begin{axis}[omaxis,
  width=9.2cm, height=5.8cm,
  xlabel={$t$}, ylabel={$T$},
  xmin=0, xmax=10.5, ymin=0, ymax=88,
  xtick={5,10}, ytick={20,50,80}]
  \addplot[omProp, dashed, domain=0:10.4] {20};
  \addplot[omDef, thick, domain=0:10.2, samples=60] {20 + 60*exp(-0.5*x)};
  \addplot[omDef!50, thick, domain=0:10.2, samples=60] {20 + 30*exp(-0.5*x)};
  \addplot[omDef!50, thick, domain=0:10.2, samples=60] {20 - 15*exp(-0.5*x)};
  \node[omProp, font=\small, anchor=south west] at (axis cs:6.6,20.8)
    {साम्यावस्था};
\end{axis}
\end{tikzpicture}

{\small आरंभिक ताप चाहे जो हो, $T' = -k(T - 20)$ के सभी हल चरघातांकी रूप से साम्य $T = 20$ पर
\omterm{def:g12:seq:limit}{अभिसरित} होते हैं।}
\end{omfigure}

\section{समीकरण \texorpdfstring{$y' = ay + f(x)$}{y' = ay + f(x)}}

\begin{proposition}[हल-समुच्चय की संरचना]\label{prop:g12:diffeq:general}
मान लीजिए $f$ किसी \omterm{def:g10:numbers:interval}{अंतराल} $I$ पर \omterm{def:g12:limcont:continuity}{संतत} है और $y_p$ $I$ पर
\[
y' = ay + f(x)
\]
का कोई एक विशेष हल है। तब $I$ पर हल ठीक \omterm{def:g11:func:function}{फलन} $y = C\eu^{ax} + y_p$, $C \in \R$ हैं।
\end{proposition}

\begin{proof}
\cref{thm:g12:diffeq:affine} की तरह: $y$ हल है यदि और केवल यदि $z = y - y_p$ $z' = (ay + f) - (ay_p + f) = az$ को संतुष्ट करे, अर्थात्
यदि और केवल यदि $z = C\eu^{ax}$ हो।
\end{proof}

\begin{method}[विशेष हल का अनुमान लगाना]\label{met:g12:diffeq:particular}
ऐसा विशेष हल ढूँढ़िए जिसका \emph{रूप $f$ जैसा} हो:
\begin{itemize}
  \item $f$ घात $n$ का बहुपद: घात $n$ का बहुपद आज़माइए;
  \item $k \neq a$ के साथ $f(x) = \alpha\,\eu^{kx}$: $y_p = \beta\,\eu^{kx}$ आज़माइए;
  \item $f(x) = \alpha\,\eu^{ax}$ (अनुनाद वाली स्थिति): $y_p = \beta x\,\eu^{ax}$ आज़माइए।
\end{itemize}
\omterm{def:g10:algebra:equation}{समीकरण} में रखिए और गुणांक मिलाइए।
\end{method}

\begin{example}\label{ex:g12:diffeq:particular}
$y' = 2y + 4x$ हल कीजिए। $y_p = \alpha x + \beta$ आज़माइए: सभी $x$ के लिए $\alpha = 2(\alpha x + \beta) + 4x$ होने से $2\alpha + 4 = 0$ और $\alpha = 2\beta$
आवश्यक हो जाते हैं, इसलिए $\alpha = -2$, $\beta = -1$। व्यापक हल: $y = C\eu^{2x} - 2x - 1$।
\end{example}

\section{अभ्यास}

\begin{exercise}[$\star$]\label{exo:g12:diffeq:1}
$\R$ पर हल कीजिए:
(a) $y(0) = 2$ के साथ $y' = 3y$;\quad
(b) $y(0) = -1$ के साथ $2y' + y = 0$;\quad
(c) $y(0) = 0$ के साथ $y' = -y + 5$।
\end{exercise}

\begin{exercise}[$\star$]\label{exo:g12:diffeq:2}
कोई जीवाणु-समष्टि अपने आकार के समानुपाती दर से बढ़ती है और हर $3$ घंटे में
दुगुनी हो जाती है। समष्टि $N(t)$ द्वारा संतुष्ट \omterm{def:g12:diffeq:ode}{अवकल समीकरण} लिखिए और
समानुपात-अचर ज्ञात कीजिए।
\end{exercise}

\begin{exercise}[$\star$]\label{exo:g12:diffeq:3}
जाँचिए कि $y_p(x) = x\,\eu^{x}$ $y' = y + \eu^x$ का एक हल है, और $\R$ पर उसके सभी हल बताइए।
\end{exercise}

\begin{exercise}[$\star\star$]\label{exo:g12:diffeq:4}
$y(0) = 1$ के साथ $y' = -2y + \eu^{x}$ हल कीजिए। (संकेत: $\beta\,\eu^{x}$ रूप का विशेष हल ढूँढ़िए।)
\end{exercise}

\begin{exercise}[$\star\star$]\label{exo:g12:diffeq:5}
कार्बन-14 $5730$ वर्ष की अर्धायु के साथ क्षय होता है। किसी पुरातात्त्विक
\omterm{def:g10:proba:sample}{प्रतिदर्श} में किसी जीवित जीव के कार्बन-14 का $60\%$ है। उसकी आयु आँकिए।
\end{exercise}

\begin{exercise}[$\star\star$]\label{exo:g12:diffeq:6}
किसी टंकी में $100$ L शुद्ध जल है। प्रति लीटर $0.2$ kg नमक वाला लवण-जल $5$
L/min से भीतर आता है, और (एकसमान बना रहने वाला) मिश्रण उसी दर से बाहर जाता है।
मान लीजिए $m(t)$ समय $t$ (मिनट में) पर टंकी में नमक का द्रव्यमान है।
\begin{enumerate}
  \item पुष्ट कीजिए कि $m' = 1 - \dfrac{m}{20}$।
  \item हल कीजिए, और $t \to +\infty$ पर $m(t)$ की सीमा ज्ञात कीजिए। उसकी व्याख्या
        कीजिए।
\end{enumerate}
\end{exercise}

\begin{exercise}[$\star\star$]\label{exo:g12:diffeq:7}
$80$ kg द्रव्यमान की एक आकाश-गोताखोर गुरुत्व ($g = 9.8\ \text{m/s}^2$) और चाल के समानुपाती
वायु-प्रतिरोध के अधीन गिरती है, जिससे उसकी चाल $k = 16$ kg/s के साथ $v' = g - \frac{k}{m}v$ को
संतुष्ट करती है।
\begin{enumerate}
  \item $v(0) = 0$ के साथ \omterm{def:g10:algebra:equation}{समीकरण} हल कीजिए।
  \item \emph{सीमांत वेग} $\lim_{t\to+\infty} v(t)$ निकालिए, और उसका $95\%$ तक पहुँचने में लगने
        वाला समय भी।
\end{enumerate}
\end{exercise}

\begin{exercise}[$\star\star\star$]\label{exo:g12:diffeq:8}
\emph{(लॉजिस्टिक \omterm{def:g10:algebra:equation}{समीकरण}.)} कोई समष्टि $y(t) \in \intoo{0}{1}$ (\omterm{def:g10:functions:extrema}{अधिकतम} समष्टि के अंश के रूप में)
\[
y' = y(1 - y).
\]
को संतुष्ट करती है।
\begin{enumerate}
  \item मान लीजिए $z = \dfrac{1}{y}$ है। दिखाइए कि $z$ रैखिक \omterm{def:g10:algebra:equation}{समीकरण} $z' = -z + 1$ को संतुष्ट
        करता है।
  \item $y(0) = \frac{1}{10}$ के साथ $z$ के लिए, फिर $y$ के लिए हल कीजिए।
  \item दिखाइए कि $t \to +\infty$ पर $y(t) \to 1$, और हल-वक्र की आकृति का रेखाचित्र बनाइए।
\end{enumerate}
\end{exercise}

\begin{exercise}[$\star\star\star$]\label{exo:g12:diffeq:9}
मान लीजिए $y$ $y' = ay + b$ का कोई हल है और $n \in \N$ के लिए $u_n = y(n)$ है। दिखाइए कि $(u_n)$
समांतर-गुणोत्तर पुनरावृत्ति $u_{n+1} = q u_n + r$ को संतुष्ट करता है, और $q$ तथा $r$ को
$a$ और $b$ के पदों में व्यक्त कीजिए। \omterm{def:g12:diffeq:ode}{अवकल समीकरण} के लिए शर्त $\abs{q} < 1$ किसके
संगत है?
\end{exercise}

\section{समस्या: वस्तुओं के भीतर की घड़ी}

\begin{problem}[{सप्ताहांत समस्या --- कार्बन-14 गुफ़ाओं की तारीख़ बताता है,
न्यूटन का शीतलन नियम अपराध का समय तय करता है, और एक छोटा-सा समीकरण चार वेश
धारण करता है}]
\label{pb:g12:diffeq:1}
\omterm{def:g12:diffeq:ode}{अवकल समीकरण} परिवर्तन का नियम है; उसे हल कर देना उस नियम को एक \emph{घड़ी} में
बदल देता है। वही नन्हा \omterm{def:g10:algebra:equation}{समीकरण} $y' = ay + b$ (\cref{thm:g12:diffeq:affine}) प्रागैतिहासिक कोयले, ठंडी होती
कॉफ़ी, गिरती बूँदों और अस्पताल की ड्रिप के भीतर टिक-टिक करता है --- और उन
घड़ियों को पढ़ना ही इस समस्या का काम है: वह लास्को की चित्रकारी की तारीख़
बताती है, मृत्यु का समय तय करती है, और एक अच्छे वैज्ञानिक की तरह अपनी ही
मान्यताओं की जाँच करती है।

\medskip
\noindent\textbf{भाग I --- प्रवाह।}
\begin{enumerate}
  \item $y' = 3y$, $y(0) = 2$ हल कीजिए; फिर $y' = -2y + 6$, $y(0) = 0$ (\cref{met:g12:diffeq:solve})।
  \item \cref{thm:g12:diffeq:homogeneous} के पीछे की अद्वितीयता वाली तरकीब फिर से निकालिए: यदि $y' = ay$ हो,
        तो $y(t)\,\eu^{-at}$ का \omterm{def:g12:deriv:derivative}{अवकलज} निकालिए और निष्कर्ष निकालिए कि \emph{हर} हल $C\eu^{at}$
        है।
  \item $y' = -2y + 6$ के लिए: साम्य हल ज्ञात कीजिए, और $t \to +\infty$ पर हर दूसरे हल की नियति
        बताइए। (\cref{pb:g11:seq:1} की पुनरावृत्तियों के \omterm{pb:g10:functions:1}{स्थिर बिंदु}, अब \omterm{def:g12:limcont:continuity}{संतत} रूप में।)
  \item क्षय के नियम $y' = -\frac{y}{\tau}$ के लिए: दिखाइए कि अर्धायु $t_{1/2} = \tau \ln 2$ है।
  \item $y' = -2y + 6$ के हलों के परिवार का रेखाचित्र बनाइए (या वर्णन कीजिए): $3$
        से ऊपर से शुरू होने वाले हल क्या करते हैं? $3$ से नीचे वाले? और
        ठीक $3$ पर वाले?
\end{enumerate}

\medskip
\noindent\textbf{भाग II --- कार्बन-14।}
जीवित ऊतक रेडियोधर्मी कार्बन-14 का अनुपात अचर बनाए रखता है; मृत्यु पर उसका
आना रुक जाता है और भंडार क्षय होने लगता है: $N' = -\lambda N$, जिसकी अर्धायु $5\,730$ वर्ष
है।
\begin{enumerate}[resume]
  \item $\lambda$ निकालिए (प्रति वर्ष)।
  \item किसी हड्डी में जीवित अनुपात का $20\,\%$ बचा है: वह कितनी पुरानी है?
  \item लास्को की चित्रित गुफ़ाओं के कोयले में लगभग $15\,\%$ बचा है: चित्रकारी
        की तारीख़ बताइए।
  \item आल्प्स के हिमनद में मिला हिम-मानव ओत्ज़ी लगभग $53\,\%$ दिखाता है: उसकी
        तारीख़ बताइए (पुरातत्त्वविद् लगभग $5\,300$ वर्ष कहते हैं --- आपका उत्तर
        कैसा रहा?)।
  \item कार्बन-14 डायनासोर की तारीख़ क्यों नहीं बता सकता? दस अर्धायुओं के
        बाद बचा हुआ अंश निकालिए, $65$ करोड़ वर्ष बाद का अंश $2$ की घात के
        रूप में बताइए, और निष्कर्ष निकालिए।
  \item त्रुटि-सीमाएँ: यदि ओत्ज़ी का $53\,\%$ केवल $\pm 1\,\%$ तक ज्ञात हो, तो आयु का
        परिसर निकालिए। प्रयोगशाला की कितनी यथार्थता तिथि-निर्धारण की कितनी
        यथार्थता ख़रीदती है?
\end{enumerate}

\medskip
\noindent\textbf{भाग III --- न्यूटन का शीतलन, और एक अपराध।}
अचर $T_a$ वाले परिवेश में ताप $T$ वाली कोई वस्तु $T' = -k\,(T - T_a)$ के अनुसार ठंडी होती
है।
\begin{enumerate}[resume]
  \item \omterm{def:g10:algebra:equation}{समीकरण} हल कीजिए ($z = T - T_a$ प्रतिस्थापित कीजिए):
        \[
        T(t) = T_a + (T_0 - T_a)\,\eu^{-kt} .
        \]
  \item $20\,^\circ$C के कमरे में $90\,^\circ$C पर उड़ेली गई कॉफ़ी $5$ मिनट बाद $70\,^\circ$C
        दिखाती है। $k$ ज्ञात कीजिए, फिर पीने योग्य $55\,^\circ$C तक की प्रतीक्षा
        का समय।
  \item दूध वाला प्रश्न: दस मिनट में सबसे गर्म कॉफ़ी पीने के लिए ठंडा दूध
        अभी डालना चाहिए या अंतिम क्षण में? उत्तर \omterm{def:g10:algebra:equation}{समीकरण} से दीजिए (दूध डालने
        से अंतर $T - T_a$ का क्या होता है, और वह अंतर हानि को कैसे चलाता है?)।
  \item न्यायालयिक विज्ञान: आधी रात को $20\,^\circ$C के कमरे में $30\,^\circ$C पर एक शव
        मिलता है; एक घंटे बाद वह $28\,^\circ$C दिखाता है। मृत्यु के समय $37\,^\circ$C
        मानते हुए दोनों मापों से $k$ निकालिए, फिर मृत्यु का समय।
  \item मृत्यु-परीक्षक की घड़ी की जाँच कीजिए: प्रतिरूप की मान्यताओं के तीन
        वास्तविक उल्लंघन बताइए और यह भी कि हर एक अनुमानित मृत्यु-समय को किस
        दिशा में झुका देगा।
\end{enumerate}

\medskip
\noindent\textbf{भाग IV --- चार वेश और एक S-वक्र।}
\begin{enumerate}[resume]
  \item \cref{exo:g12:diffeq:8} का लॉजिस्टिक \omterm{def:g10:algebra:equation}{समीकरण}, $y(0) = 0.1$ के साथ $y' = y(1 - y)$: उसी अभ्यास वाले
        प्रतिस्थापन से $y(t) = \dfrac{1}{1 + 9\eu^{-t}}$ निकालिए, और \omterm{def:g12:deriv:inflection}{नति-परिवर्तन बिंदु} का समय $y = \frac12$
        ज्ञात कीजिए --- महामारी-विज्ञानी की \cref{pb:g12:deriv:1} वाली तारीख़, अब गणनीय।
  \item कोई वर्षा-बूँद $v' = 10 - \frac v2$ (गुरुत्व घटा कर्षण) मानती है, $v(0) = 0$ के साथ।
        सीमांत वेग ज्ञात कीजिए, $v(t)$ के लिए हल कीजिए, और निकालिए कि बूँद
        सीमांत चाल का $95\,\%$ कब पा लेती है।
  \item अस्पताल की ड्रिप अचर दर से दवा देती है जबकि शरीर उसे समानुपाती रूप
        से निकालता जाता है: $c' = 4 - 0.5c$, $c(0) = 0$। स्थायी-अवस्था की सांद्रता और उसका
        आधा पाने का समय ज्ञात कीजिए। \cref{exo:g12:diffeq:9} की ऋण-पुनरावृत्तियों से समांतरता
        एक वाक्य में बताइए।
  \item समापन --- एक \omterm{def:g10:algebra:equation}{समीकरण}, चार वेश: क्षय, शीतलन, गिरना, ख़ुराक --- सब
        अलग-अलग चिह्नों और नामों वाले $y' = ay + b$। इस समस्या ने जो प्रतिरूपण-चक्र
        चार बार चलाया उसे दोहराइए (नियम $\to$ \omterm{def:g10:algebra:equation}{समीकरण} $\to$ हल $\to$
        अंशांकन $\to$ भविष्यवाणी $\to$ जाँच), और उस परिघटना का नाम बताइए
        जिसे $y''$ चाहिए --- और वह भी कि कौन-सी समस्या उससे पहले ही मिल चुकी
        है।
\end{enumerate}
\end{problem}
